\begin{abstract}
Emerging advancements exemplified by BitNet~\cite{bitnet} are ushering in a paradigm shift toward 1-bit Large Language Models (LLMs). In this paper, we present a ternary LLM variant, \textbf{\text{BitNet b1.58}}, whose parameters are restricted to the set \{-1, 0, 1\}. This architecture achieves parity with conventional full-precision Transformers (employing FP16 or BF16) of equivalent scale and exposure to training tokens in terms of both perplexity and downstream task outcomes. Simultaneously, it yields dramatic improvements in efficiency—including reductions in latency, memory utilization, throughput demands, and energy expenditure. Crucially, the introduction of 1.58-bit LLMs establishes a novel scaling law and training paradigm, facilitating the development of subsequent LLM generations that reconcile performance excellence with resource efficiency. Additionally, this direction enables entirely new computational strategies and lays the groundwork for bespoke hardware designs optimized for 1-bit LLM inference and training.
\end{abstract}

\section{The Era of 1-bit LLMs}

The last several years have witnessed the explosive expansion of Large Language Models (LLMs) within AI, fueling unprecedented advances across diverse natural language processing benchmarks. However, the growing scale of these systems brings with it substantial obstacles for practical deployment, as well as increased environmental and financial burdens resulting from significant energy demands.

To counteract these issues, researchers have adopted post-training quantization strategies, compressing model parameters to a reduced bit-width for inference~\cite{smoothquant, gptq, quip, quip_sharp}. By lowering the numerical precision of both weights and activations, these approaches markedly cut down the computational and memory costs of LLMs. The prevailing trajectory in quantization research has moved from 16 bits to even more compact representations—such as 4-bit schemes~\cite{gptq, awq}. Despite the popularity and commercial adoption of post-training quantization, these methods remain inherently sub-optimal compared to other quantization and compression paradigms.

Advancements in 1-bit model designs, exemplified by BitNet~\cite{bitnet}, demonstrate significant potential for lowering large language model (LLM) costs without compromising efficacy. Standard LLM implementations utilize 16-bit floating-point formats (such as FP16 or BF16), and most computational overhead arises from extensive matrix multiplications, which predominantly consist of floating-point addition and multiplication. Conversely, BitNet streamlines this process by reducing matrix multiplication to only integer addition, yielding substantial energy savings for LLM inference and training. Since energy consumption is a primary performance bottleneck for many hardware chips, reducing it not only cuts costs but also enables accelerated computation.

Beyond computation, model parameter transfer from DRAM to on-chip accelerator memory (e.g., SRAM) during inference is another significant expense. Efforts to increase SRAM size to boost throughput are often hampered by its high cost relative to DRAM. Compared to traditional full-precision models, 1-bit LLMs substantially decrease both memory usage and bandwidth requirements. This reduction dramatically lessens the cost and time associated with moving weights from DRAM, facilitating faster and more efficient inference.

In this paper, we present a notable extension to 1-bit LLMs called \textbf{\text{BitNet b1.58}}, where model parameters are ternary, adopting values in \{-1, 0, 1\}. By introducing the 0 value into BitNet's original 1-bit formulation, this architecture achieves a bit-width of 1.58 bits in binary encoding. \text{BitNet b1.58} preserves all of the original's benefits, including an efficient computational framework with nearly multiplication-free matrix multiplications, readily amenable to hardware optimization. Its energy consumption remains on par with classic 1-bit BitNet models, while offering significantly improved memory efficiency, throughput, and latency compared to FP16-based LLMs. Moreover, \text{BitNet b1.58} provides two notable advantages: enhanced expressiveness from direct feature filtering enabled by zero weights—which markedly bolsters modeling capacity—and empirical evidence that, starting with models of 3B parameters or larger under matched conditions (e.g., model architecture, number of training tokens), \text{BitNet b1.58} can match full-precision (FP16) counterparts on both perplexity and downstream tasks.

\section{BitNet b1.58}

\text{BitNet b1.58} extends the BitNet framework, a variant of the Transformer in which the conventional \emph{nn.Linear} layers are substituted with \emph{BitLinear} layers. This model is trained entirely anew, employing 1.58-bit quantized weights and 8-bit activations. The architecture features several adjustments relative to the original BitNet, which are outlined below.

\paragraph{Quantization Function.} We utilize an \emph{absmean} quantization strategy to enforce that weights take on values exclusively in \{-1, 0, +1\}. This process involves normalizing the weight matrix by its mean absolute value before rounding each entry to the nearest element in the target set:

\begin{equation}
    \widetilde{W} = \mathrm{RoundClip}(\frac{W}{\gamma + \epsilon}, -1, 1),
\end{equation}

\begin{equation}
    \mathrm{RoundClip}(x, a, b) = \max(a, \min(b, \mathrm{round}(x))),
\end{equation}

\begin{equation}
    \gamma = \frac{1}{nm}\sum_{ij} |W_{ij}|.
\end{equation}

For activation quantization, we adhere to the approach used in BitNet but introduce a key difference: activations are not scaled to fit within $[0, Q_b]$ before applying non-linearities. Rather, activations for each token are normalized to fall within $[-Q_b, Q_b]$, thereby eliminating the necessity for zero-point quantization. This alternative is not only easier to implement and facilitates system-level optimizations but also has an insignificant impact on empirical performance according to our results.

\paragraph{LLaMA-alike Components.}

Given that LLaMA~\cite{llama, llama2} has become the standard template for open-source large language models, we adopt its architectural elements to align \text{BitNet b1.58} with prevailing open-source practices. Concretely, our model incorporates RMSNorm~\cite{rmsnorm}, SwiGLU~\cite{swiglu}, and rotary positional embeddings~\cite{rope}, while omitting all biases. This alignment ensures straightforward compatibility with widely used open-source tools (such as Huggingface, vLLM~\cite{vllm}, and llama.cpp\footnote{\url{https://github.com/ggerganov/llama.cpp}}), greatly reducing integration overhead.

\section{Results}

We conducted a comprehensive comparison between \text{BitNet b1.58} and a re-implemented FP16 version of \text{LLaMA LLM} across a spectrum of model scales. To guarantee an equitable evaluation, both models underwent pre-training for 100 billion tokens using the RedPajama dataset~\cite{redpajama}. Zero-shot task performance was assessed on several benchmark language understanding tasks, such as ARC-Easy~\cite{arc}, ARC-Challenge~\cite{arc}, Hellaswag~\cite{hellaswag}, Winogrande~\cite{winoGrande}, PIQA~\cite{piqa}, OpenbookQA~\cite{openbookqa}, and BoolQ~\cite{boolq}. Furthermore, validation perplexity was reported for WikiText2~\cite{wikitext2} and C4~\cite{c4} datasets.

In addition, we measured and compared GPU memory consumption and inference latency for \text{LLaMA LLM} and \text{BitNet b1.58}. All performance metrics were obtained via the FasterTransformer\footnote{\url{https://github.com/NVIDIA/FasterTransformer}} framework, optimized for efficient LLM inference on GPUs. For \text{BitNet b1.58}, the Ladder~\cite{ladder} 2-bit kernel was integrated into the evaluation. Latency was quantified as the time per generated output token, reflecting the principal inference cost.

Table~\ref{tab:ppl} presents a summary of perplexity metrics alongside computational cost for both \text{BitNet b1.58} and \text{LLaMA LLM}. The findings indicate that at the 3B parameter scale, \text{BitNet b1.58} attains parity in perplexity with its full-precision \text{LLaMA LLM} counterpart, while delivering a 2.71-fold increase in speed and reducing GPU memory usage by a factor of 3.55. Notably, \text{BitNet b1.58} at 3.9B parameters achieves 2.4x faster inference and consumes 3.32x less GPU memory, yet still surpasses the 3B \text{LLaMA LLM} in downstream performance.

A detailed breakdown of zero-shot accuracy across downstream tasks can be found in Table~\ref{tab:zero-shot}. We employed the evaluation setup from \emph{lm-evaluation-harness}\footnote{\url{https://github.com/EleutherAI/lm-evaluation-harness}} for these experiments. The results reveal that the accuracy difference between \text{BitNet b1.58} and \text{LLaMA LLM} diminishes as model size grows. Crucially, starting from the 3B parameter scale, \text{BitNet b1.58} achieves performance equivalent to full-precision baselines. These trends in perplexity are mirrored by task performance; specifically, the 3.9B \text{BitNet b1.58} outstrips the 3B \text{LLaMA LLM} in both accuracy and resource efficiency. Altogether, this substantiates \text{BitNet b1.58} as a Pareto-optimal enhancement over current leading LLMs.

\paragraph{Memory and Latency}

We expanded our analysis to encompass model sizes of 7B, 13B, and 70B, and conducted a cost evaluation. As shown in Figure~\ref{fig:latency-memory-energy}, both latency and memory usage demonstrate clear scaling behaviors, with greater speed-up realized as models become larger. Notably, at the 70B scale, \text{BitNet b1.58} achieves a speed advantage, running 4.1 times faster than the \text{LLaMA LLM} reference. This acceleration arises primarily because the computational burden associated with \emph{nn.Linear} operations grows with increasing model scale. Memory usage follows a parallel pattern: as models grow, the share of memory devoted to embeddings (which are kept in full precision) declines proportionally. All latency and memory benchmarks utilized a 2-bit kernel, indicating that further optimizations are possible to decrease costs even more.

\paragraph{Energy}

Our assessment also includes an estimation of the energy required for arithmetic operations in both \text{BitNet b1.58} and \text{LLaMA LLM}. Since matrix multiplication dominates large language model computational costs, our evaluation centers on this operation. Figure~\ref{fig:energy} breaks down the components contributing to total energy usage: for \text{BitNet b1.58}, the majority consists of INT8 addition operations, whereas \text{LLaMA LLM} expends energy through both FP16 addition and multiplication. Drawing on the energy models from~\cite{energycost, pokebnn}, we find that \text{BitNet b1.58} reduces the energy spent on arithmetic operations for matrix multiplication by a factor of 71.4 on 7nm hardware. We also provide comprehensive end-to-end energy measurements for runs involving 512 tokens. These results indicate that as model size grows, \text{BitNet b1.58} becomes progressively more energy-efficient relative to the FP16 \text{LLaMA LLM} baseline. This is largely attributable to the increasing proportion of time and energy spent on \emph{nn.Linear} in larger models, while contributions from other subsystems become less significant.

\paragraph{Throughput}

We evaluate the throughput of \text{BitNet b1.58} compared to the \text{LLaMA LLM} with 70B parameters, both deployed across two 80GB A100 GPUs using pipeline parallelism~\cite{gpipe}—enabling \text{LLaMA LLM} 70B to run within the available hardware constraints. The batch size was scaled up incrementally until GPU memory was saturated, maintaining a fixed sequence length of 512. As shown in Table~\ref{tab:throughput}, \text{BitNet b1.58} 70B is capable of accommodating a batch size up to 11 times larger than \text{LLaMA LLM}, resulting in an 8.9-fold increase in throughput.

\textbf{\text{BitNet b1.58} demonstrates a novel scaling relationship between model size, computational performance, and inference efficiency}. To illustrate, equivalences between 1.58-bit and 16-bit model sizes, as observed from Figures~\ref{fig:latency-memory-energy} and \ref{fig:energy}, can be described as follows:

\begin{itemize}
    \item 13B BitNet b1.58 surpasses the 3B FP16 LLM regarding latency, memory footprint, and energy demand.
    \item 30B BitNet b1.58 delivers superior latency, memory consumption, and energy efficiency compared to the 7B FP16 LLM.
    \item 70B BitNet b1.58 outperforms the 13B FP16 LLM in latency, memory, and power usage.
\end{itemize}

\paragraph{Training with 2T Tokens}

The total number of tokens used during training significantly influences LLM performance. To assess the token scalability of \text{BitNet b1.58}, we trained this model on 2T tokens using the data preparation method from StableLM-3B~\cite{StableLM-3B-4E1T}, which represents a leading open-source 3B model. Both models underwent evaluation on a test suite comprising Winogrande~\cite{winoGrande}, PIQA~\cite{piqa}, SciQ~\cite{sciq}, LAMBADA~\cite{lambada}, and ARC-easy~\cite{arc}. Zero-shot accuracy scores, summarized in Table~\ref{tab:2t}, average results across tasks assessed via both accuracy and normalized accuracy. For StableLM 3B at 2T tokens, outcomes are drawn directly from its technical report. Our analysis indicates that \text{BitNet b1.58} consistently achieves higher scores on every task, suggesting robust generalization capabilities for 1.58-bit LLMs.

\section{Discussion and Future Work}

\textbf{1-bit Mixture-of-Experts (MoE) LLMs}

Mixture-of-Experts (MoE) architectures have established themselves as a resource-efficient solution for scaling LLMs. Although these models decrease computational FLOPs, their deployment is often constrained by substantial memory demands and high inter-device communication costs. Transitioning to 1.58-bit LLMs can alleviate these issues. The drastic reduction in memory requirements allows for deploying MoE models using fewer devices, and inter-chip communication for activation transfer is significantly diminished. In the ideal scenario, if the entire model fits onto a single chip, inter-device overhead is eliminated entirely.

\textbf{Native Support of Long Sequence in LLMs}

As LLMs progress, accommodating longer sequences has emerged as a vital requirement. A primary obstacle for efficient long sequence inference is the extensive memory usage associated with KV caches. \text{BitNet b1.58} marks notable progress towards inherent long-sequence support by halving activation precision from 16 bits to 8 bits, effectively doubling the possible context window within the same memory constraints. Additionally, further lossless compression to 4 bits—or possibly less with 1.58-bit LLMs—could be achieved, which we propose as a direction for future exploration.

\textbf{LLMs on Edge and Mobile}

Deploying 1.58-bit LLMs on edge and mobile platforms can notably extend the reach and functionality of language models. The stringent memory and computational limitations of such devices have traditionally hindered the adoption and scalability of LLMs. The drastically lower energy and memory requirements of 1.58-bit LLMs make their deployment feasible on resource-constrained hardware, unlocking new application domains previously out of reach. As a result, the abilities of edge and mobile platforms are greatly enhanced, opening opportunities for innovative LLM applications. Additionally, 1.58-bit LLMs are particularly suited to CPUs, the predominant processing units in these environments. This compatibility ensures that \text{BitNet b1.58} can operate with high efficiency on such devices, further expanding their utility.

\textbf{New Hardware for 1-bit LLMs}

Recent developments such as~Groq\footnote{\url{https://groq.com/}} have illustrated the significant advantages of dedicated hardware platforms (for instance, LPUs) tailored for LLM workloads. Building on these efforts, we advocate for designing novel hardware and systems expressly optimized for 1-bit LLMs, taking full advantage of the computational efficiencies introduced by BitNet~\cite{bitnet}.